\DTMsavedate{mydate}{2022-06-30}
\maketitle{Microprocessors}{Electronic Engineering}{\DTMusedate{mydate}}

% ----------------------------------------------------- %
% COURSE INFO
\subsection*{Course Information}\label{course:tee2112}
\vspace{0ex}%
\noindent\parbox{0.5\textwidth}{%
    \noindent\begin{tabular}{@{}ll}
                 \textsf{Course:}          & 	Microprocessors    \\
                 \textsf{Course Code:}         & 	TEE 2112    \\
                 \textsf{Credit Hours:}    & 	10
    \end{tabular}}
\vspace{2ex}\hrule\vspace{2ex}

% ----------------------------------------------------- %
% INSTRUCTOR INFO
\subsection*{Lecturer Information}
\noindent\parbox{0.5\textwidth}{%
    \noindent\begin{tabular}{@{}ll}
                 \textsf{Name:}         &     Reginald Gonye         \\
%\textsf{Phone:}        &    \phone          \\
\textsf{Email:}        &    reginald.gonye@nust.ac.zw \\
\textsf{Qualifications:}        &  MPhil in Electronic Engineering, (Hon) BEng in Electronic Engineering
    \end{tabular}}
\vspace{2ex}\hrule\vspace{2ex}


% ----------------------------------------------------- %
% COURSE DESCRIPTION
\subsection*{Learning Aim}
To introduce students to microprocessors and microprocessor based systems (computers).
\vspace{2ex} \hrule\vspace{2ex}

% ----------------------------------------------------- %
% LEARNING OBJECTIVES
\subsection*{Learning Objectives}
By the end of the course, candidates should be able to:
\begin{outline}
    \1      Understand what a Microprocessor is and how it works.

    \1  	Appreciate the necessary components and understand the required hardware connections between the components that make up a microprocessor-based system or “computer” as it is known.

    \1  	Understand microprocessor instructions and their representation using assembly language statements.

    \1  	Design programs to solve novel situations using microprocessors.

    \1  	Appreciate the operation of microprocessor-based systems.

\end{outline}
\vspace{2ex}\hrule\vspace{2ex}

% ----------------------------------------------------- %
% COURSE TOPICS
\subsection*{Topics}
\begin{outline}
    \1 Introduction
    \1 Data Representation
    \1 Main Memory for Microprocessor
    \1 Input/Output for a Microprocessor
    \1 Microprocessor Internal Architecture
    \1 Microprocessor Instructions
    \1 Designing Microprocessor Programs
    \1 Stack, Subroutines and Interrupts
    \1 Control of Peripherals
    \1 Internal Operation of Microprocessors
    \1 Application of Microprocessors and Performance
\end{outline}
\vspace{2ex} \hrule\vspace{2ex}

% ----------------------------------------------------- %
% Students Assenssment
\subsection*{Assessment of Students}
\begin{outline}
    \1 Continuous assessment will consist of tests, assignments and laboratory work.
    \1 Final assessment will consist of an examination at the end of the semester. The examination shall comprise of three sections as follows:
    \begin{enumerate}[label=\alph*)]
        \item SECTION A – answer all questions; mainly factual recall; total of 40 marks
        \item SECTION B – 3 questions each with 15 marks; answer any two; questions will examine at a level of candidates’ understanding of a microprocessor (understand what programs do, transforming basic operations into programs, signal sequences as instructions execute, etc); total of 30 marks
        \item SECTION C – 3 questions each with 30 marks; answer any one; questions will examine the candidate’s ability to design a microprocessor-based solution to a novel situation; total of 30 marks.
    \end{enumerate}
\end{outline}
\vspace{2ex}\hrule\vspace{2ex}

% ----------------------------------------------------- %
% Grade Evaluation
\subsection*{Course Evaluation and Grading Policies}
Grades are based on:
Continuous assessment (\qty{25}{\percent}),
Final Exam (\qty{75}{\percent})
\vspace{2ex}\hrule\vspace{2ex}

% ----------------------------------------------------- %
% EQUIPMENT
% https://sites.udel.edu/computing-purchases/personal-specs/
% \subsection*{Essential Equipment}

% \vspace{2ex}\hrule\vspace{2ex}

% Policies and Other information
\subsection*{Policies and Other Information}
% Conduct
\subsection*{Key Text}

\begin{itemize}
    \item Microprocessor Theory and Applications with 68000/68020 and Pentium, 2008 by R. Mohamed
    \item Microprocessors and Microprocessor-based System Design, 2nd Edition, 1995 by R. Mohamed
    \item Introduction to Microprocessors, 3rd Edition, 1989 by A. P. Mathur
    \item The Z80 manual
    \item Microprocessors: Essentials, Components and Systems by R. Meadows, A. J. Parsons
\end{itemize}
\vspace{2ex}\hrule
\vspace{2ex}\hrule\vspace{2ex}